\documentclass[11pt]{article}

\usepackage[a4paper,margin=1in]{geometry}
\usepackage{amsmath,amssymb}
\usepackage{graphicx}
\usepackage{xcolor}
\usepackage{booktabs}
\usepackage{float}
\usepackage{caption}
\usepackage{hyperref}

\newcommand{\up}{\uparrow}
\newcommand{\down}{\downarrow}
\newcommand{\ee}{e}
\newcommand{\Nj}{N_J}
\newcommand{\Ne}{N_e}
\newcommand{\GammaC}{\Gamma}
\newcommand{\OmegaDrive}{\Omega}

\setlength{\parindent}{0pt}

\title{Mean-Field Equations for Inhomogeneous Couplings}
\author{}
\date{}

\begin{document}
\maketitle


\section{Inhomogeneous Hamiltonian}

Split the atoms into two groups according to the drive they feel. The Hamiltonian changes as
\begin{equation}
\hat{H}_{\delta}
=
\Omega \hat{J}_x
-
\delta \hat{N}_e
\quad
\rightarrow
\quad
\hat{H}_{\delta}
=
\Omega
(
\omega_1 \hat{J}_{1x}
+
\omega_2 \hat{J}_{2x}
)
-
\delta
(
\hat{N}_{e1}
+
\hat{N}_{e2}
).
\end{equation}

The natural Dicke/occupation basis is now described by the number of excitations in each subgroup,
\begin{equation}
|n_e\rangle
\quad
\rightarrow
\quad
|n_{e1},n_{e2}\rangle
=
|n_{e1}\rangle_1
\otimes
|n_{e2}\rangle_2.
\end{equation}

The Schwinger-boson operators have been split for the two groups,
\begin{equation}
\hat{e},\hat{d}
\quad
\rightarrow
\quad
\hat{e}_1,\hat{e}_2,\hat{d}_1,\hat{d}_2.
\end{equation}
They act as before, but only on their own subgroup. They follow the bosonic commutation relations
\begin{equation} \label{eq:bosonic_commutation_relations}
\begin{gathered}
[\hat{d}_a,\hat{d}_b^{\dagger}]
=
[\hat{e}_a,\hat{e}_b^{\dagger}]
=
\delta_{ab}.
\\
[\hat{d}_a,\hat{e}_b]
=
[\hat{d}_a,\hat{e}_b^{\dagger}]
=
[\hat{d}_a^{\dagger},\hat{e}_b]
=
[\hat{d}_a^{\dagger},\hat{e}_b^{\dagger}]
=
0.
\end{gathered}
\end{equation}
\\
The operators in the Hamiltonian act on their respective subspaces. For the excitation number operator,
\begin{equation}
\hat{N}_e
=
\sum_i |e_i\rangle\langle e_i|
\quad
\rightarrow
\quad
\hat{N}_{e1},
\hat{N}_{e2},
\end{equation}
where
\begin{equation}
\hat{N}_{e1}
=
\sum_{i\in 1}
|e_i\rangle\langle e_i|,
\qquad
\hat{N}_{e2}
=
\sum_{i\in 2}
|e_i\rangle\langle e_i|.
\end{equation}

In Schwinger-boson notation,
\begin{equation}
\hat{N}_e
=
\hat{e}^{\dagger}\hat{e}
\quad
\rightarrow
\quad
\hat{N}_e
=
\hat{N}_{e1}
+
\hat{N}_{e2}
=
\hat{e}_1^{\dagger}\hat{e}_1
+
\hat{e}_2^{\dagger}\hat{e}_2.
\end{equation}

Here the group operators act only on their own subgroup. For example,
\begin{equation}
\hat{e}_1 |n_{e1},n_{e2}\rangle
=
\sqrt{n_{e1}}
|n_{e1}-1,n_{e2}\rangle,
\end{equation}
while acting as the identity on the second subgroup.
\\

For \(\hat{J}_x\),
\begin{equation}
\hat{J}_x
=
\frac{1}{2}
(
\hat{J}_{+}
+
\hat{J}_{-}
)
=
\frac{1}{2}
(
\hat{e}^{\dagger}\hat{d}
+
\hat{d}^{\dagger}\hat{e}
).
\end{equation}

After splitting into two groups,
\begin{equation}
\hat{J}_x
=
\hat{J}_{1x}
+
\hat{J}_{2x},
\end{equation}
with
\begin{equation}
\hat{J}_{1x}
=
\frac{1}{2}
(
\hat{J}_{1+}
+
\hat{J}_{1-}
),
\qquad
\hat{J}_{2x}
=
\frac{1}{2}
(
\hat{J}_{2+}
+
\hat{J}_{2-}
).
\end{equation}

The raising and lowering operators are
\begin{equation} 
\begin{aligned}
\hat{J}_{1+}
=
\hat{e}_1^{\dagger}\hat{d}_1,
\qquad
\hat{J}_{1-}
=
\hat{d}_1^{\dagger}\hat{e}_1,
\\
\hat{J}_{2+}
=
\hat{e}_2^{\dagger}\hat{d}_2,
\qquad
\hat{J}_{2-}
=
\hat{d}_2^{\dagger}\hat{e}_2.
\end{aligned}
\end{equation}



\section{Mean-field approximation}

Replace operator expectation values by complex amplitudes,
\begin{equation}
D_a(t)
\equiv
\langle \hat{d}_a\rangle(t),
\qquad
E_a(t)
\equiv
\langle \hat{e}_a\rangle(t),
\end{equation}
for \(a=1,2\).
\\

Approximate
\begin{equation}
|D_a|^2 
=
|\langle \hat{d}_a\rangle|^2 = \langle \hat{d}_a\rangle \langle \hat{d}_a\rangle 
\approx
\langle \hat{d}_a^{\dagger}\hat{d}_a\rangle
=
N_{d,a},
\end{equation}
and
\begin{equation}
|E_a|^2
=
|\langle \hat{e}_a\rangle|^2 = \langle \hat{e}_a\rangle \langle \hat{e}_a\rangle
\approx
\langle \hat{e}_a^{\dagger}\hat{e}_a\rangle
=
N_{e,a}.
\end{equation}

For a coherent spin state in group \(a\),
\begin{equation}
D_a
=
\sqrt{N_{J,a}}\cos\frac{\theta_{J,a}}{2},
\qquad
E_a
=
\sqrt{N_{J,a}}e^{-i\phi_{J,a}}\sin\frac{\theta_{J,a}}{2}.
\end{equation}

More generally, the two groups may have different steady-state angles and phases. Note that the amplitudes $D_a$ and $E_a$ are defined within a two-level system with fixed $N_J$, the $J$ index has been omitted from the amplitude notation.

\section{Equations of motion}

In the Heisenberg picture, the Lindblad equation gives
\begin{equation}
\frac{d}{dt}\langle \hat{O}\rangle
=
i\langle [\hat{H},\hat{O}]\rangle
+
\Gamma
\left\langle
\hat{L}^{\dagger}\hat{O}\hat{L}
-
\frac{1}{2}
\left\{
\hat{L}^{\dagger}\hat{L},\hat{O}
\right\}
\right\rangle,
\end{equation}
with weighted collective decay
\begin{equation}
\hat{L}
=
\omega_1 \hat{J}_{1-}
+
\omega_2 \hat{J}_{2-}.
\end{equation}
Applying the bosonic commutation relations in Eq.~\ref{eq:bosonic_commutation_relations}, one finds the following mean-field equations for $\hat d_1$, $\hat d_2$, $\hat e_1$ and $\hat e_2$:
\begin{equation}
\begin{cases}
\partial_t D_a
=
-\dfrac{i\Omega\omega_a}{2}E_a
+
\dfrac{\Gamma \omega_a }{2}
(
\omega_1 E_1^{*}D_1
+
\omega_2 E_2^{*}D_2
)
E_a,
\\[2mm]
\partial_t E_a
=
-\dfrac{i\Omega\omega_a}{2}D_a
+
i\delta E_a
-
\dfrac{\Gamma \omega_a }{2}
D_a
(
\omega_1 D_1^{*}E_1
+
\omega_2 D_2^{*}E_2
),
\end{cases}
\qquad
a=1,2.
\end{equation}

% These are the four mean-field equations for the two groups.

\section{Rotating frame}

Now go to a rotating frame,
\begin{equation}
D_a(t)
=
e^{-i\omega_{B,a} t}d_a(t),
\qquad
E_a(t)
=
e^{-i\omega_{B,a} t}e_a(t).
\end{equation}

Then
\begin{equation}
\partial_t D_a
=
e^{-i\omega_{B,a} t}
(
\partial_t d_a
-
i\omega_{B,a} d_a
),
\end{equation}
and similarly for \(E_a\). Hence,
\begin{equation}
\begin{gathered}
\partial_t d_a
=
-\frac{i\Omega\omega_a}{2}e_a
+
\frac{\Gamma \omega_a}{2}
(
\omega_1 e_1^{*}d_1
+
\omega_2 e_2^{*}d_2
)
e_a
+
i\omega_{B,a} d_a,
\\
\partial_t e_a
=
-\frac{i\Omega\omega_a}{2}d_a
+
i\delta e_a
-
\frac{\Gamma \omega_a}{2}
d_a
(
\omega_1 d_1^{*}e_1
+
\omega_2 d_2^{*}e_2
)
+
i\omega_{B,a} e_a.
\end{gathered}
\end{equation}

Set
\begin{equation}
\partial_t d_a=0,
\qquad
\partial_t e_a=0.
\end{equation}

Multiply the first equation by \(d_a^{*}\), and multiply the complex conjugate of the second equation by \(e_a\). This gives
\begin{equation}
0
=
-\frac{i\Omega\omega_a}{2}e_a d_a^{*}
+
\frac{\Gamma \omega_a}{2}
(
\omega_1 e_1^{*}d_1
+
\omega_2 e_2^{*}d_2
)
e_a d_a^{*}
+
i\omega_{B,a} |d_a|^2,
\end{equation}
and
\begin{equation}
0
=
\frac{i\Omega\omega_a}{2}d_a^{*}e_a
-
i\delta |e_a|^2
-
\frac{\Gamma \omega_a}{2}
d_a^{*}e_a
(
\omega_1 d_1e_1^{*}
+
\omega_2 d_2e_2^{*}
)
-
i\omega_{B,a} |e_a|^2.
\end{equation}

Adding the two equations cancels the drive terms and dissipative terms:
\begin{equation}
0
=
-i\delta |e_a|^2
+
i\omega_{B,a}
(
|d_a|^2-|e_a|^2
).
\end{equation}

Using
\begin{equation}
|d_a|^2
=
N_{J,a}\cos^2\frac{\theta_{J,a}}{2},
\qquad
|e_a|^2
=
N_{J,a}\sin^2\frac{\theta_{J,a}}{2},
\end{equation}
we find
\begin{equation} \label{eq:rotating_frame}
\omega_{B,a}
=
\frac{
\delta 
}{
2\cos\theta_{J,a}
} - \frac{\delta}{2}.
\end{equation}
If we enforce that the two groups share the same rotating frame $\omega_{B,1} = \omega_{B,2}$, the above equation implies that they share the same excitation ratio as well then, meaning  $\theta_1 = \theta_2$. Therefore the inhomogeneous drive weights \(\omega_1\) and \(\omega_2\) must enter mainly through different phases \(\phi_1\) and \(\phi_2\) in that case.
\\

Now subtract the second equation from the first. This gives
\begin{equation}
0
=
-i\Omega\omega_a d_a^{*}e_a
+
i\delta |e_a|^2
+
i\omega_{B,a}
(
|d_a|^2+|e_a|^2
)
+
\Gamma \omega_a d_a^{*}e_a
(
\omega_1 d_1e_1^{*}
+
\omega_2 d_2e_2^{*}
).
\end{equation}

After substituting
\begin{equation}
\begin{gathered}
d_a^{*}e_a
=
\frac{N_{J,a}}{2}
e^{-i\phi_{J,a}}
\sin\theta_{J,a},
\\
\omega_1 d_1e_1^{*}
+
\omega_2 d_2e_2^{*}
=
\frac{\omega_1 N_{J,1}}{2}
e^{i\phi_{J,1}}
\sin\theta_{J,1}
+
\frac{\omega_2 N_{J,2}}{2}
e^{i\phi_{J,2}}
\sin\theta_{J,2},
\end{gathered}
\end{equation}
and using the formula found for $\omega_B$ in Eq. \ref{eq:rotating_frame}, one arrives at:
\begin{equation} \label{eq:mfe}
\boxed{
\begin{aligned}
\frac{\Omega\omega_a}{2}
e^{-i\phi_{J,a}}
\sin\theta_{J,a}
=
&\frac{\delta}{2}
\sin \theta_{J,a} \tan \theta_{J,a}
-
i\frac{\Gamma \omega_a}{4}
e^{-i\phi_{J,a}}
\sin\theta_{J,a}
\left[
\omega_1 N_{J,1} e^{i\phi_{J,1}} \sin\theta_{J,1}
+
\omega_2 N_{J,2}e^{i\phi_{J,2}} \sin\theta_{J,2}
\right].
\end{aligned}
}
\end{equation}
This is the steady-state phase equation for group \(a=1,2\).

\section{Homogeneous limit}
In the homogeneous limit of Eq. \ref{eq:mfe},
\begin{equation}
\omega_1=\omega_2=1,
\qquad
\theta_{J,1}=\theta_{J,2}=\theta_J,
\qquad
\phi_{J,1}=\phi_{J,2}=\phi_J.
\end{equation}

Then
\begin{equation}
N_{J,1}e^{i\phi_{J,1}}
+
N_{J,2}e^{i\phi_{J,2}}
=
N_J e^{i\phi_J},
\end{equation}
and the group equation reduces to
\begin{equation}
\begin{aligned}
\frac{\Omega}{2}
e^{-i\phi_J}
\sin\theta_J
=
&\frac{\delta}{2}
\sin \theta_J \tan \theta_J
-
i\frac{\Gamma N_J}{4}
\sin^2\theta_J.,
\end{aligned}
\end{equation}
matching perfectly the equation for the homogeneous case derived in the original paper. 

\section{Wavefunction coefficients (fix)}
In homogeneous case, one single active manifold defined by  ....
\\
...
\\

Two-group sector coefficients given by
\begin{equation}
P(N_{J,1},N_{J,2})
=
\binom{N_1}{N_{J,1}}
\binom{N_2}{N_{J,2}}
2^{-N}.
\end{equation}


\section{Critical drive}

Taking the real part of the mean field equations in Eq. \ref{eq:mfe}, and simplifying gives

\begin{equation}
\begin{aligned}
\frac{\Omega}{2}
=
-\frac{\Gamma}{4}
\left[
\omega_1 N_{J,1} \sin \phi_{J,1} \sin\theta_{J,1}
+
\omega_2 N_{J,2} \sin \phi_{J,2} \sin\theta_{J,2}
\right].
\end{aligned}
\end{equation}
In the central $(N_{J,1}, N_{J,2}) = (N_1/2,N_2/2)$ case
\begin{equation}
\begin{aligned}
\frac{\Omega}{2}
& \leq
\frac{\Gamma}{4}
\left[
\omega_1 N_{J,1}
+
\omega_2 N_{J,2} 
\right] =
\\
& = \frac{\Gamma}{4} \frac{1}{2} \left[ \omega_1 N_1 + \omega_2 N_2 \right]
\\
& = \frac{\Gamma}{4} \frac{N}{2}
\\
& = \frac{\Gamma N_J}{4}
\end{aligned}
\end{equation}
Therefore, the same critical drive may be assumed in the inhomogeneous setting as the homogeneous one:
\begin{equation}
    \Omega_c = \frac{\Gamma N_J}{2}
\end{equation}
where $N_J = N/2$, if effective couplings satisfy $\omega_1 N_1 + \omega_2 N_2 = N$.

\end{document}